%% Auto-generated completion supplement for AcademicWriting.tex and TASK_FROM_SCIENTIFIC_ADVISOR.txt
%% Source files: /botworkspace/THESIS/TEX/AcademicWriting.tex (lines 1191--1213),
%%               /botworkspace/THESIS/TASK_FROM_SCIENTIFIC_ADVISOR.txt
%% Purpose: Complete formulations and mathematical proofs for the "WANTED BUT FORGOT",
%%          "WANTED BUT TOO EARLY", and Scientific Advisor tasks.

\section{Omitted Definitions and Advisor Tasks}
\label{sec:wanted_but_forgot_completions}

\subsection{Weak Measurability, Weak Topologies, and Mixed Norms}
\label{subsec:dunford_pettis_prerequisites}

Let $(M, \Sigma, \mu)$ be a $\sigma$-finite measure space and let $E$ be a Banach space over $\mathbb{K} \in \{\mathbb{R}, \mathbb{C}\}$ with topological dual $E'$. We denote the canonical duality pairing between $\varphi \in E'$ and $f \in E$ by $\langle \varphi, f \rangle = \varphi(f)$.

\begin{definition}[Weakly measurable function unique up to sets of measure zero]
\label{def:weakly_measurable_unique}
A function $a : M \to E'$ is called \emph{weak-$*$ measurable} (or \emph{$E$-weakly measurable} in the terminology of Simon~\cite{Simon1982} and Pettis~\cite{pettis1938integration}) if for every $f \in E$, the scalar-valued function
\[
x \mapsto \langle a(x), f \rangle
\]
is $\Sigma$-measurable on $M$. When $E = L^p(M, d\mu)$ with $1 \le p < \infty$ and $E' = L^{p'}(M, d\mu)$ ($p' = p/(p-1)$), two weak-$*$ measurable functions $a_1, a_2 : M \to E'$ are said to be \emph{equivalent up to sets of measure zero} if there exists a single $\mu$-null set $N \in \Sigma$ (with $\mu(N) = 0$) independent of $f \in E$ such that $a_1(x) = a_2(x)$ in $E'$ for all $x \in M \setminus N$.
\end{definition}

\begin{lemma}[Uniqueness of weak-$*$ measurable representatives on separable pre-duals]
\label{lem:weak_measurable_uniqueness}
Assume that the Banach space $E$ is separable (which holds in particular for $E = L^p(M, d\mu)$ when $1 \le p < \infty$ and $(M, \Sigma, \mu)$ is countably generated and $\sigma$-finite, such as $M = \mathbb{R}^\nu$ with Lebesgue measure). Let $a_1, a_2 : M \to E'$ be two weak-$*$ measurable functions such that for every $f \in E$,
\[
\langle a_1(x), f \rangle = \langle a_2(x), f \rangle \quad \text{for $\mu$-almost every } x \in M.
\]
Then there exists a single $\mu$-null set $N \in \Sigma$ with $\mu(N) = 0$ such that $a_1(x) = a_2(x)$ in $E'$ for all $x \in M \setminus N$. Moreover, the norm function $x \mapsto \|a_1(x)\|_{E'}$ is $\Sigma$-measurable, and
\[
\operatorname*{ess\,sup}_{x \in M} \|a_1(x)\|_{E'} = \sup_{\|f\|_E \le 1} \|\langle a_1(\cdot), f \rangle\|_{L^\infty(M, d\mu)}.
\]
\end{lemma}

\begin{proof}
Since $E$ is separable, choose a countable dense subset $D = \{f_n : n \in \mathbb{N}\} \subset E$. For each $n \in \mathbb{N}$, by hypothesis there exists a measurable set $N_n \in \Sigma$ with $\mu(N_n) = 0$ such that
\[
\langle a_1(x), f_n \rangle = \langle a_2(x), f_n \rangle \quad \text{for all } x \in M \setminus N_n.
\]
Define $N = \bigcup_{n=1}^\infty N_n \in \Sigma$. By countable subadditivity of $\mu$, we have $\mu(N) \le \sum_{n=1}^\infty \mu(N_n) = 0$, so $\mu(N) = 0$. Fix any $x \in M \setminus N$. Then the continuous linear functionals $a_1(x), a_2(x) \in E'$ agree on the dense subset $D \subset E$. Since both $a_1(x)$ and $a_2(x)$ are norm-continuous on $E$ and coincide on $D$, they coincide on all of $E$; that is, $a_1(x) = a_2(x)$ in $E'$ for every $x \in M \setminus N$.

Furthermore, if $D_1 = \{g_n : n \in \mathbb{N}\}$ is a countable dense subset of the closed unit ball $B_E = \{f \in E : \|f\|_E \le 1\}$, then for every $x \in M$,
\[
\|a_1(x)\|_{E'} = \sup_{f \in B_E} |\langle a_1(x), f \rangle| = \sup_{n \in \mathbb{N}} |\langle a_1(x), g_n \rangle|.
\]
Being the pointwise supremum of a countable family of measurable nonnegative functions, $x \mapsto \|a_1(x)\|_{E'}$ is $\Sigma$-measurable, and taking essential suprema commutes with the countable supremum over $n \in \mathbb{N}$, yielding the stated identity.
\end{proof}

\begin{definition}[Weak topology $\sigma(C, \mathcal{M})$]
\label{def:weak_topology_sigma_C_M}
Let $C$ be a vector space over $\mathbb{K}$ and let $\mathcal{M}$ be a vector space of linear functionals on $C$ (or, more generally, a space in duality with $C$ via a bilinear pairing $\langle \cdot, \cdot \rangle : \mathcal{M} \times C \to \mathbb{K}$). The \emph{weak topology on $C$ induced by $\mathcal{M}$}, denoted by $\sigma(C, \mathcal{M})$ (following Tr\`eves~\cite{Treves1967}), is the coarsest (weakest) topology on $C$ with respect to which every linear functional $\ell \in \mathcal{M}$ is continuous. Equivalently, $\sigma(C, \mathcal{M})$ is the locally convex vector topology generated by the family of seminorms
\[
p_F(u) = \max_{\ell \in F} |\langle \ell, u \rangle|, \quad u \in C,
\]
as $F$ ranges over all non-empty finite subsets of $\mathcal{M}$. A net $(u_\alpha)$ in $C$ converges to $u \in C$ in $\sigma(C, \mathcal{M})$ if and only if $\langle \ell, u_\alpha \rangle \to \langle \ell, u \rangle$ in $\mathbb{K}$ for every $\ell \in \mathcal{M}$. The topology $\sigma(C, \mathcal{M})$ is Hausdorff if and only if $\mathcal{M}$ separates points of $C$ (that is, $\langle \ell, u \rangle = 0$ for all $\ell \in \mathcal{M}$ implies $u = 0$).
\end{definition}

\begin{definition}[Mixed operator norm $\| A \|_{p;\infty}$]
\label{def:operator_norm_p_infty}
Let $(M, \Sigma, \mu)$ be a measure space and let $1 \le p \le \infty$. For any bounded linear operator $A : L^p(M, d\mu) \to L^\infty(M, d\mu)$, the \emph{$(p, \infty)$-operator norm} $\|A\|_{p;\infty}$ is defined by
\[
\|A\|_{p;\infty} := \|A\|_{\mathcal{B}(L^p(M, d\mu),\, L^\infty(M, d\mu))}
= \sup_{\substack{f \in L^p(M, d\mu) \\ \|f\|_{L^p} \le 1}} \|Af\|_{L^\infty(M, d\mu)}.
\]
By the Dunford--Pettis theorem (Simon~\cite[Theorem A.1.1]{Simon1982}), when $1 \le p < \infty$ and $L^p(M, d\mu)$ is separable, every bounded linear operator $A : L^p(M, d\mu) \to L^\infty(M, d\mu)$ is an integral operator represented by a unique (up to $\mu$-null sets) weak-$*$ measurable kernel $a : M \to L^{p'}(M, d\mu)$ ($1/p + 1/p' = 1$) such that for each $f \in L^p(M, d\mu)$,
\[
(Af)(x) = \int_M a(x)(y)\, f(y)\, d\mu(y) \quad \text{for $\mu$-almost every } x \in M,
\]
and the operator norm satisfies the isometric identity
\[
\|A\|_{p;\infty} = \operatorname*{ess\,sup}_{x \in M} \|a(x)\|_{L^{p'}(M, d\mu)}.
\]
\end{definition}

\subsection{Schr\"odinger Operators and Semigroups}
\label{subsec:schroedinger_operators_completed}

\begin{definition}[Schr\"odinger operator]
\label{def:schroedinger_operator_completed}
On the Hilbert space $\mathcal{H} = L^2(\mathbb{R}^\nu)$ ($\nu \ge 1$), the \emph{free Schr\"odinger operator} (or kinetic energy operator) is the non-negative self-adjoint operator
\[
H_0 := -\frac{1}{2} \Delta = -\frac{1}{2} \sum_{j=1}^{\nu} \frac{\partial^2}{\partial x_j^2},
\]
with maximal Sobolev domain
\[
D(H_0) = H^2(\mathbb{R}^\nu) = \left\{ \psi \in L^2(\mathbb{R}^\nu) : |\xi|^2 \widehat{\psi}(\xi) \in L^2(\mathbb{R}^\nu) \right\},
\]
where $\Delta$ operates in the sense of tempered distributions (and Note: in the PDE convention used in Chapter 1, one writes $H_0^{\mathrm{PDE}} = -\Delta = 2 H_0$, which differs only by the time rescaling $t \mapsto 2t$). Given a real-valued measurable potential $V : \mathbb{R}^\nu \to \mathbb{R}$, let $M_V$ (denoted simply by $V$) be the multiplication operator $(V\psi)(x) = V(x)\psi(x)$ on $D(V) = \{\psi \in L^2(\mathbb{R}^\nu) : V\psi \in L^2(\mathbb{R}^\nu)\}$. The \emph{Schr\"odinger operator} is the operator sum
\[
H = H_0 + V.
\]
\end{definition}

\begin{proposition}[Self-adjointness and Schr\"odinger semigroup for bounded potentials]
\label{prop:schroedinger_semigroup_bounded}
If $V \in L^\infty(\mathbb{R}^\nu; \mathbb{R})$, then $M_V$ is a bounded self-adjoint operator on $L^2(\mathbb{R}^\nu)$ with $\|M_V\|_{\mathcal{B}(L^2)} = \|V\|_{L^\infty}$. Consequently, by the Kato--Rellich theorem, $H = H_0 + V$ is self-adjoint on $D(H) = D(H_0) = H^2(\mathbb{R}^\nu)$ and is bounded from below by $-\|V\|_{L^\infty}$:
\[
\langle H\psi, \psi \rangle = \frac{1}{2} \|\nabla \psi\|_{L^2}^2 + \int_{\mathbb{R}^\nu} V(x) |\psi(x)|^2\, dx \ge -\|V\|_{L^\infty} \|\psi\|_{L^2}^2 \quad \text{for all } \psi \in H^2(\mathbb{R}^\nu).
\]
By the spectral theorem for self-adjoint operators, $H$ generates a strongly continuous holomorphic semigroup of bounded self-adjoint operators $\{\exp(-tH)\}_{t \ge 0}$ on $L^2(\mathbb{R}^\nu)$ satisfying the operator norm bound $\|\exp(-tH)\|_{\mathcal{B}(L^2)} \le \exp(t \|V\|_{L^\infty})$.
\end{proposition}

\subsection{Solutions to Tasks from the Scientific Advisor}
\label{subsec:advisor_tasks_solutions}

\subsubsection{Advisor Task 1: Verification of Kernel Assumptions for \texorpdfstring{$\exp(-tH)$}{exp(-tH)}}
\label{subsubsec:advisor_task_1}

\begin{proposition}[Kernel and semigroup assumptions for $\exp(-tH_0)$ and $\exp(-tH)$]
\label{prop:advisor_task_1_heat_kernel}
Let $H_0 = -\frac{1}{2}\Delta$ on $L^2(\mathbb{R}^\nu)$ (or $H_0^{\mathrm{PDE}} = -\Delta$). For every $t > 0$, the operator $\exp(-tH_0)$ is an integral operator on $L^2(\mathbb{R}^\nu)$ with explicit Gaussian heat kernel
\[
p_t(x, y) = \frac{1}{(2\pi t)^{\nu/2}} \exp\left( -\frac{|x - y|^2}{2t} \right), \quad x, y \in \mathbb{R}^\nu
\]
(or $p_t^{\mathrm{PDE}}(x, y) = (4\pi t)^{-\nu/2} \exp(-|x - y|^2 / (4t))$ for $\exp(t\Delta)$). This family satisfies all abstract kernel assumptions:
\begin{enumerate}
\item \textbf{Positivity and normalization:} $p_t(x, y) > 0$ for all $x, y \in \mathbb{R}^\nu$, $p_t(x, y) = p_t(y, x)$, and $\int_{\mathbb{R}^\nu} p_t(x, y)\, dy = 1$ for all $x \in \mathbb{R}^\nu$.
\item \textbf{Ultracontractivity ($L^p$ to $L^\infty$ boundedness):} For every $1 \le p \le \infty$ and $t > 0$, $\exp(-tH_0)$ maps $L^p(\mathbb{R}^\nu)$ continuously into $L^\infty(\mathbb{R}^\nu)$, with mixed norm
\[
\|\exp(-tH_0)\|_{p;\infty} = \sup_{x \in \mathbb{R}^\nu} \|p_t(x, \cdot)\|_{L^{p'}(\mathbb{R}^\nu)} = (2\pi t)^{-\frac{\nu}{2p}} (p')^{-\frac{\nu}{2p'}} < \infty,
\]
where $1/p + 1/p' = 1$ (with $(p')^{-1/p'} = 1$ for $p' = \infty$). In particular, by the Dunford--Pettis theorem, $\exp(-tH_0)$ possesses a bounded integral kernel.
\item \textbf{Semigroup homomorphism (Chapman--Kolmogorov identity):} The map $t \mapsto \exp(-tH_0)$ is a semigroup homomorphism from $((0, \infty), +)$ into $\mathcal{B}(L^2(\mathbb{R}^\nu))$, corresponding to the convolution identity
\[
\int_{\mathbb{R}^\nu} p_s(x, z)\, p_t(z, y)\, dz = p_{s+t}(x, y) \quad \text{for all } s, t > 0 \text{ and } x, y \in \mathbb{R}^\nu.
\]
\item \textbf{Translation covariance in the spatial variable:} Because $p_t(x, y) = k_t(x - y)$ depends only on the difference $x - y$, for every shift $a \in \mathbb{R}^\nu$ we have $p_t(x - a, y - a) = p_t(x, y)$, so $\exp(-tH_0)$ commutes with the unitary group of spatial translations $\{T_a\}_{a \in \mathbb{R}^\nu}$.
\item \textbf{Perturbed Schr\"odinger semigroup $\exp(-tH)$:} If $V \in L^\infty(\mathbb{R}^\nu; \mathbb{R})$, then by the Feynman--Kac formula (or Dyson Phillips perturbation series), $\exp(-tH)$ is an integral operator with positive symmetric kernel $K_t(x, y)$ satisfying the pointwise Gaussian bounds
\[
\exp(-t\|V\|_{L^\infty})\, p_t(x, y) \le K_t(x, y) \le \exp(t\|V\|_{L^\infty})\, p_t(x, y),
\]
so $\|\exp(-tH)\|_{p;\infty} \le \exp(t\|V\|_{L^\infty}) \|\exp(-tH_0)\|_{p;\infty} < \infty$ for all $1 \le p \le \infty$ and $t > 0$.
\end{enumerate}
\end{proposition}

\begin{proof}
Under the unitary Fourier transform $(\mathcal{F}f)(\xi) = (2\pi)^{-\nu/2} \int_{\mathbb{R}^\nu} e^{-ix\cdot \xi} f(x)\, dx$, the operator $H_0 = -\frac{1}{2}\Delta$ is unitarily equivalent to multiplication by $\frac{1}{2}|\xi|^2$. By functional calculus,
\[
(\mathcal{F}(\exp(-tH_0)f))(\xi) = e^{-t|\xi|^2/2} (\mathcal{F}f)(\xi).
\]
Since the inverse Fourier transform of $(2\pi)^{-\nu/2} e^{-t|\xi|^2/2}$ is the Gaussian kernel
\[
k_t(x) = (2\pi t)^{-\nu/2} e^{-|x|^2/(2t)},
\]
the convolution theorem yields
\[
(\exp(-tH_0)f)(x) = (k_t * f)(x) = \int_{\mathbb{R}^\nu} p_t(x, y) f(y)\, dy.
\]

For (2), Young's convolution inequality (or H\"older's inequality) implies that $\|\exp(-tH_0)f\|_{L^\infty} \le \|k_t\|_{L^{p'}} \|f\|_{L^p}$, and direct integration of the Gaussian gives
\[
\|k_t\|_{L^{p'}}^{p'} = (2\pi t)^{-\frac{\nu p'}{2}} \int_{\mathbb{R}^\nu} \exp\left(-\frac{p'|z|^2}{2t}\right) dz
= (2\pi t)^{-\frac{\nu p'}{2}} \left(\frac{2\pi t}{p'}\right)^{\frac{\nu}{2}} = (2\pi t)^{-\frac{\nu(p'-1)}{2}} (p')^{-\frac{\nu}{2}}.
\]
Taking the $1/p'$-th power and noting $(p'-1)/p' = 1/p$ yields the exact value of $\|\exp(-tH_0)\|_{p;\infty}$. Property (3) follows immediately from $e^{-s|\xi|^2/2} e^{-t|\xi|^2/2} = e^{-(s+t)|\xi|^2/2}$ in Fourier space, and property (5) follows from Trotter's product formula
\[
\exp(-t(H_0+V)) = \operatorname*{s-lim}_{n\to\infty} \bigl(\exp(-(t/n)H_0)\exp(-(t/n)V)\bigr)^n
\]
together with $|\exp(-(t/n)V(x))| \le \exp((t/n)\|V\|_{L^\infty})$.
\end{proof}

\subsubsection{Advisor Task 2: Multiplication Operator Spectrum, Essential Range, Fourier Conjugation, and Dense Testing}
\label{subsubsec:advisor_task_2}

\begin{definition}[Essential range]
\label{def:essential_range}
Let $(X, \Sigma, \mu)$ be a $\sigma$-finite measure space and let $m : X \to \mathbb{C}$ be a measurable function. The \emph{essential range} of $m$, denoted $\operatorname{ess\,ran}(m)$, is the set
\[
\operatorname{ess\,ran}(m) := \left\{ \lambda \in \mathbb{C} : \forall \varepsilon > 0,\ \mu\big(\{x \in X : |m(x) - \lambda| < \varepsilon\}\big) > 0 \right\}.
\]
\end{definition}

\begin{theorem}[Spectrum of a multiplication operator equals the essential range]
\label{thm:multiplication_operator_spectrum}
Let $(X, \Sigma, \mu)$ be a $\sigma$-finite measure space and let $m : X \to \mathbb{C}$ be measurable. Define the multiplication operator $M_m$ on $\mathcal{H} = L^2(X, d\mu)$ by
\[
D(M_m) = \{\psi \in L^2(X, d\mu) : m\psi \in L^2(X, d\mu)\}, \quad (M_m \psi)(x) = m(x)\psi(x).
\]
Then $M_m$ is a densely defined closed normal operator (it is self-adjoint if and only if $\operatorname{ess\,ran}(m) \subset \mathbb{R}$, and bounded if and only if $m \in L^\infty(X, d\mu)$, with $\|M_m\| = \|m\|_{L^\infty}$), and its spectrum satisfies
\[
\sigma(M_m) = \operatorname{ess\,ran}(m).
\]
\end{theorem}

\begin{proof}
\emph{Step 1: $\mathbb{C} \setminus \operatorname{ess\,ran}(m) \subset \rho(M_m)$.}
Suppose $\lambda \notin \operatorname{ess\,ran}(m)$. By definition, there exists $\varepsilon > 0$ such that the set $E_\varepsilon := \{x \in X : |m(x) - \lambda| < \varepsilon\}$ has measure $\mu(E_\varepsilon) = 0$. Consequently, $|m(x) - \lambda| \ge \varepsilon$ for $\mu$-almost every $x \in X$, so the function $r_\lambda(x) := (m(x) - \lambda)^{-1}$ belongs to $L^\infty(X, d\mu)$ with $\|r_\lambda\|_{L^\infty} \le \varepsilon^{-1}$. The bounded multiplication operator $M_{r_\lambda} \in \mathcal{B}(L^2(X, d\mu))$ satisfies $(M_m - \lambda I) M_{r_\lambda} \psi = \psi$ for all $\psi \in L^2(X, d\mu)$ and $M_{r_\lambda} (M_m - \lambda I)\psi = \psi$ for all $\psi \in D(M_m)$. Thus $M_m - \lambda I$ has a bounded two-sided inverse, so $\lambda \in \rho(M_m)$. Taking complements gives $\sigma(M_m) \subset \operatorname{ess\,ran}(m)$.

\emph{Step 2: $\operatorname{ess\,ran}(m) \subset \sigma(M_m)$.}
Suppose $\lambda \in \operatorname{ess\,ran}(m)$. Then for every $n \in \mathbb{N}$, the measurable set
\[
A_n := \left\{ x \in X : |m(x) - \lambda| < \frac{1}{n} \right\}
\]
has $\mu(A_n) > 0$. Since $(X, \Sigma, \mu)$ is $\sigma$-finite, we can choose a measurable subset $B_n \subset A_n$ such that $0 < \mu(B_n) < \infty$. Define the unit vector $\psi_n := \mu(B_n)^{-1/2} \mathbf{1}_{B_n} \in D(M_m) \subset L^2(X, d\mu)$, which satisfies $\|\psi_n\|_{L^2} = 1$ and $|m(x)| \le |\lambda| + 1$ on $B_n$. Then
\[
\|(M_m - \lambda I)\psi_n\|_{L^2}^2 = \frac{1}{\mu(B_n)} \int_{B_n} |m(x) - \lambda|^2\, d\mu(x) \le \frac{1}{n^2}.
\]
Thus $\|(M_m - \lambda I)\psi_n\|_{L^2} \le 1/n \to 0$ as $n \to \infty$ while $\|\psi_n\|_{L^2} = 1$ (so $(\psi_n)$ is a Weyl sequence for $\lambda$). Therefore $M_m - \lambda I$ cannot have a bounded inverse on $L^2(X, d\mu)$, meaning $\lambda \in \sigma(M_m)$. Combining both inclusions yields $\sigma(M_m) = \operatorname{ess\,ran}(m)$.
\end{proof}

\begin{lemma}[Almost-sure equality from a countable dense family of test vectors]
\label{lem:dense_testing_almost_sure}
Let $(X, \Sigma, \mu)$ be a $\sigma$-finite measure space such that $L^2(X, d\mu)$ is separable, and let $\{\psi_n\}_{n \in \mathbb{N}}$ be a countable dense subset of $L^2(X, d\mu)$ (or of $L^1(X, d\mu) \cap L^2(X, d\mu)$).
\begin{enumerate}
\item If $f, g \in L^2(X, d\mu)$ satisfy $\langle f, \psi_n \rangle_{L^2} = \langle g, \psi_n \rangle_{L^2}$ for all $n \in \mathbb{N}$, then $f(x) = g(x)$ for $\mu$-almost every $x \in X$.
\item If $A, B \in \mathcal{B}(L^2(X, d\mu))$ satisfy $(A\psi_n)(x) = (B\psi_n)(x)$ for $\mu$-almost every $x \in X$ for every $n \in \mathbb{N}$, then $A = B$ in $\mathcal{B}(L^2(X, d\mu))$, and for every $\psi \in L^2(X, d\mu)$, $(A\psi)(x) = (B\psi)(x)$ for $\mu$-almost every $x \in X$.
\end{enumerate}
\end{lemma}

\begin{proof}
For (1), put $h = f - g \in L^2(X, d\mu)$. By hypothesis, $\langle h, \psi_n \rangle = 0$ for every $n \in \mathbb{N}$. Given $h \in L^2(X, d\mu)$, since $\{\psi_n\}$ is dense in $L^2(X, d\mu)$, there exists a subsequence $(\psi_{n_k})$ converging to $h$ in $L^2(X, d\mu)$. By continuity of the inner product,
\[
\|h\|_{L^2}^2 = \langle h, h \rangle = \lim_{k \to \infty} \langle h, \psi_{n_k} \rangle = 0,
\]
so $\int_X |f(x) - g(x)|^2\, d\mu(x) = 0$, which implies $f(x) = g(x)$ for $\mu$-almost every $x \in X$.

For (2), for each $n \in \mathbb{N}$, the hypothesis $(A\psi_n)(x) = (B\psi_n)(x)$ $\mu$-almost everywhere means $\|A\psi_n - B\psi_n\|_{L^2} = 0$. Since $A - B$ is a continuous linear operator on $L^2(X, d\mu)$ that vanishes on the dense subset $\{\psi_n\}_{n \in \mathbb{N}}$, it vanishes identically on $L^2(X, d\mu)$. Hence $\|A\psi - B\psi\|_{L^2} = 0$ for every $\psi \in L^2(X, d\mu)$, which means $(A\psi)(x) = (B\psi)(x)$ for $\mu$-almost every $x \in X$.
\end{proof}

\subsubsection{Advisor Task 3: Exercise on Page 92 --- Fourier Transform of the Shift Operator}
\label{subsubsec:advisor_task_3}

\begin{proposition}[Fourier transform of the shift operator and unitary group homomorphism]
\label{prop:shift_operator_fourier}
For each translation vector $a \in \mathbb{R}^\nu$, define the shift (translation) operator $T_a : L^2(\mathbb{R}^\nu) \to L^2(\mathbb{R}^\nu)$ by
\[
(T_a f)(x) := f(x - a), \quad x \in \mathbb{R}^\nu
\]
(and similarly the left-shift $(S_a f)(x) := f(x + a) = (T_{-a} f)(x)$). Let $\mathcal{F} : L^2(\mathbb{R}^\nu) \to L^2(\mathbb{R}^\nu)$ denote the unitary Fourier--Plancherel transform defined on $f \in L^1(\mathbb{R}^\nu) \cap L^2(\mathbb{R}^\nu)$ by
\[
(\mathcal{F}f)(\xi) = \widehat{f}(\xi) := \frac{1}{(2\pi)^{\nu/2}} \int_{\mathbb{R}^\nu} e^{-i x \cdot \xi}\, f(x)\, dx, \quad \xi \in \mathbb{R}^\nu.
\]
Then:
\begin{enumerate}
\item For every $a \in \mathbb{R}^\nu$ and $f \in L^2(\mathbb{R}^\nu)$,
\[
\big(\mathcal{F}(T_a f)\big)(\xi) = e^{-i a \cdot \xi}\, (\mathcal{F}f)(\xi) \quad \text{and} \quad \big(\mathcal{F}(S_a f)\big)(\xi) = e^{i a \cdot \xi}\, (\mathcal{F}f)(\xi)
\]
for Lebesgue-almost every $\xi \in \mathbb{R}^\nu$.
\item Consequently, $\mathcal{F}$ unitarily diagonalizes $T_a$ into a multiplication operator:
\[
\mathcal{F} T_a \mathcal{F}^{-1} = M_{e_a}, \quad \text{where } e_a(\xi) := e^{-i a \cdot \xi}.
\]
\item The map $a \mapsto T_a$ is a strongly continuous unitary group representation (homomorphism) from the additive group $(\mathbb{R}^\nu, +)$ into the unitary group $\mathcal{U}(L^2(\mathbb{R}^\nu))$:
\[
T_0 = I, \quad T_a T_b = T_{a+b}, \quad T_a^* = T_a^{-1} = T_{-a} \quad \text{for all } a, b \in \mathbb{R}^\nu,
\]
and by Theorem~\ref{thm:multiplication_operator_spectrum}, for every $a \ne 0$, the spectrum of $T_a$ is the unit circle:
\[
\sigma(T_a) = \sigma(M_{e_a}) = \operatorname{ess\,ran}(e_a) = \mathbb{T} = \{z \in \mathbb{C} : |z| = 1\}.
\]
\end{enumerate}
\end{proposition}

\begin{proof}
First take $f \in \mathcal{S}(\mathbb{R}^\nu) \subset L^1(\mathbb{R}^\nu) \cap L^2(\mathbb{R}^\nu)$. By the translation invariance of Lebesgue measure under the substitution $y = x - a$ ($x = y + a$, $dx = dy$), we compute directly:
\begin{align*}
\big(\mathcal{F}(T_a f)\big)(\xi)
&= \frac{1}{(2\pi)^{\nu/2}} \int_{\mathbb{R}^\nu} e^{-i x \cdot \xi}\, f(x - a)\, dx \\
&= \frac{1}{(2\pi)^{\nu/2}} \int_{\mathbb{R}^\nu} e^{-i (y + a) \cdot \xi}\, f(y)\, dy \\
&= e^{-i a \cdot \xi} \frac{1}{(2\pi)^{\nu/2}} \int_{\mathbb{R}^\nu} e^{-i y \cdot \xi}\, f(y)\, dy
= e^{-i a \cdot \xi}\, (\mathcal{F}f)(\xi).
\end{align*}
Replacing $a$ by $-a$ yields $(\mathcal{F}(S_a f))(\xi) = e^{i a \cdot \xi} (\mathcal{F}f)(\xi)$. Since $\mathcal{S}(\mathbb{R}^\nu)$ is dense in $L^2(\mathbb{R}^\nu)$ and both $\mathcal{F} T_a$ and $M_{e_a} \mathcal{F}$ are bounded linear isometries on $L^2(\mathbb{R}^\nu)$ (as $|e^{-i a \cdot \xi}| = 1$), Lemma~\ref{lem:dense_testing_almost_sure}(2) extends the identity $\mathcal{F} T_a f = M_{e_a} \mathcal{F} f$ to all $f \in L^2(\mathbb{R}^\nu)$. Conjugating on the right by $\mathcal{F}^{-1}$ gives $\mathcal{F} T_a \mathcal{F}^{-1} = M_{e_a}$.

Finally, in Fourier space, the group homomorphism laws reduce immediately to the exponential law for characters:
\[
e_0(\xi) = 1, \quad e_a(\xi) e_b(\xi) = e^{-i a \cdot \xi} e^{-i b \cdot \xi} = e^{-i(a+b)\cdot \xi} = e_{a+b}(\xi), \quad \overline{e_a(\xi)} = e^{i a \cdot \xi} = e_{-a}(\xi) = e_a(\xi)^{-1}.
\]
Because conjugation by the unitary operator $\mathcal{F}$ is a $*$-isomorphism of $\mathcal{B}(L^2(\mathbb{R}^\nu))$, this implies $T_0 = I$, $T_a T_b = T_{a+b}$, and $T_a^* = T_{-a} = T_a^{-1}$. For $a \ne 0$, the continuous map $\xi \mapsto e^{-i a \cdot \xi}$ maps $\mathbb{R}^\nu$ onto the unit circle $\mathbb{T}$, and the preimage of every open arc of $\mathbb{T}$ has infinite Lebesgue measure, so $\sigma(T_a) = \operatorname{ess\,ran}(e_a) = \mathbb{T}$.
\end{proof}
